\section{Síntesis y ampliación}

Hyung Won Chung (钟亨原) completó con un deck de 67 páginas una «arqueología de arquitecturas»: primero planteó la fuerza direccional de un exponentially cheaper compute; después usó la Bitter Lesson para explicar la tensión entre estructura y escala; luego descompuso el paso de encoder-decoder a decoder-only en cuatro transformaciones locales, y por último recurrió a la traducción, a FLAN, a las representaciones jerárquicas y a la caché en conversaciones de varios turnos para explicar por qué cada sesgo fue razonable en su momento y por qué puede haber caducado.

\subsection{Diez conclusiones centrales}

\begin{enumerate}
  \item Al estudiar un campo que cambia con rapidez, además de seguir los resultados hay que estudiar la trayectoria del cambio.
  \item La dominant driving force es una aproximación de modelado con márgenes de error, no una declaración de causa única.
  \item La tendencia histórica de compute per dollar puede orientar la estrategia de investigación, pero no demuestra por sí sola la pendiente futura.
  \item El inductive bias aporta a la vez eficiencia presente y un posible scale ceiling.
  \item ``Scalable'' significa que el desempeño sigue mejorando a medida que crece el presupuesto; no equivale a poder ejecutarse de forma distribuida.
  \item Encoder-decoder, encoder-only y decoder-only comparten la base de sequence-model, y sus diferencias pueden descomponerse localmente.
  \item Las cuatro transformation retiran, respectivamente, los supuestos estructurales sobre los parámetros de attention, los parámetros del stack, la layer connection y la mask.
  \item La traducción automática y los datos de long-input/short-target explican por qué encoder-decoder tuvo ventaja en su momento.
  \item El valor tanto de la final-layer-only cross-attention como de la bidirectionality debe volver a medirse según la profundidad y el modo de interacción.
  \item Lo que realmente ofrece la historia de las arquitecturas es un bias audit, no una lista eterna de ganadores.
\end{enumerate}

\subsection{Una tabla que condensa el ciclo de vida de las estructuras}

\begin{table}[H]
\centering
\small
\begin{tabular}{@{}L{0.17\textwidth}L{0.24\textwidth}L{0.24\textwidth}L{0.25\textwidth}@{}}
\toprule
Supuesto estructural & Justificación en el entorno anterior & Cambio de regime & Prueba que conviene hacer \\
\midrule
Separar los parámetros de entrada y de objetivo & En la traducción, los dos lados tienen papeles y vocabularios distintos & Conocimiento general, generación larga, diálogo de varios turnos & Scaling curve con parámetros compartidos \\
Leer solo la última encoder layer & Suficiente con profundidad moderada y sencillo de implementar & Redes más profundas, jerarquías de granularidad más ricas & Layer-aligned cross-attention \\
Attention bidireccional sobre la entrada & Tareas de comprensión sobre la entrada completa & Streaming e interacción incremental de varios turnos & Ganancia de calidad frente al costo de recodificar \\
Architecture label fija & Facilita la implementación y la comparación & Componentes sustituibles localmente, cambios en las restricciones del sistema & Ablation de una sola variable por transformation \\
\bottomrule
\end{tabular}
\caption{Expresar las estructuras históricas como hipótesis de ciclo de vida que pueden verificarse.}
\end{table}

\subsection{Práctica: el procedimiento mínimo de un bias audit}

El siguiente pseudocódigo convierte el método de la clase en un registro de investigación. No sustituye a los experimentos, pero evita que el equipo reporte solo la mejor puntuación actual y olvide registrar las condiciones en las que la estructura es válida.

\begin{lstlisting}[caption={Plantilla de registro de investigación para un inductive-bias audit}]
for bias in architecture.assumptions:
    document(original_problem(bias))
    measure(current_regime_gain(bias))
    choose_scale_axes(compute, data, depth, interaction)
    run_single_variable_ablation(bias, scale_axes)
    report(crossover_point, uncertainty, system_cost)
\end{lstlisting}

\begin{lstlisting}[caption={Pseudocódigo de decisión para distinguir el óptimo actual de la tendencia de largo plazo}]
if deployment_budget < observed_crossover:
    ship(currently_efficient_method)
else:
    validate(less_structured_method)
keep_both_curves_and_failure_cases()
never_label_extrapolation_as_measurement()
\end{lstlisting}

\begin{warningbox}{No convertir la «investigación por sustracción» en un nuevo dogma}
Retirar estructura no es, por naturaleza, más científico. Si un sesgo expresa una invariancia real, reduce de forma notable la necesidad de datos o cumple restricciones de latencia, privacidad y consumo de energía, puede seguir siendo razonable a largo plazo. Lo correcto es medir beneficios, costos y el punto de cruce, no tratar ``less structure'' como una preferencia estética que no requiere evidencia.
\end{warningbox}

\subsection{Preguntas de autoevaluación}

\begin{enumerate}
  \item ¿Por qué el experimento de dejar caer una pluma permite ignorar la resistencia del aire, pero no permite concluir que un sistema complejo tiene siempre una sola variable dominante?
  \item ¿En qué se diferencian compute per dollar y la Moore's law? ¿Qué es lo que la figura de la clase no demuestra?
  \item ¿Cómo puede el inductive bias aumentar la sample efficiency y, al mismo tiempo, reducir la escalabilidad a largo plazo?
  \item ¿Cuáles son las tres attention role del encoder-decoder?
  \item ¿Qué supuesto estructural retira cada uno de los cuatro pasos de transformation?
  \item ¿Por qué la traducción automática favorece de forma natural la separación de parámetros de entrada y de objetivo?
  \item ¿Por qué la mayor ganancia de FLAN-T5 solo respalda la hipótesis sobre la distribución de los datos, en lugar de constituir una prueba causal?
  \item ¿Qué cuello de botella puede provocar la final-layer-only cross-attention en una red muy profunda?
  \item ¿Por qué la bidirectional attention impide reutilizar la KV cache en un diálogo de varios turnos?
  \item Si la tendencia del compute cambia, ¿qué supuestos deben actualizarse, en lugar de conservar mecánicamente cuál conclusión?
\end{enumerate}

\subsection{Lecturas adicionales}

\begin{itemize}
  \item Richard Sutton, \emph{The Bitter Lesson}: la relación de largo plazo entre los métodos generales, el cómputo y la estructura diseñada por personas.
  \item Vaswani et al., \emph{Attention Is All You Need}: el Transformer encoder-decoder original y los tres tipos de attention role.
  \item Devlin et al., \emph{BERT}: el diseño representativo de encoder-only y del masked language modeling.
  \item Brown et al., \emph{Language Models are Few-Shot Learners}: un hito histórico del scaling de decoder-only.
  \item Chung et al., \emph{Scaling Instruction-Finetuned Language Models}: la evidencia de instruction-tuning de FLAN-T5/PaLM.
\end{itemize}

\end{document}
